%%%%%%%%%%%%%%%%%%%%%%%%%%%%%%%%%%%%%%%%%%%%%%%%%%%%%%%%%%%%%%%%%%%%%%%%%%%%%%%%%%%%%%%%%%%%%%%%%%%%%%%%%%%%%%%%%%%%%%%
%
%  Filename   : preamble.tex
%
%  Description: Preamble file to :
%               a. specify related packages
%              F b. set margins, commands, etc.
%
%  Note       : Edit the margin settings for your own printer
%                  You may add your own commands, environments (it is assumed that you know what you're doing.)
%
%%%%%%%%%%%%%%%%%%%%%%%%%%%%%%%%%%%%%%%%%%%%%%%%%%%%%%%%%%%%%%%%%%%%%%%%%%%%%%%%%%%%%%%%%%%%%%%%%%%%%%%%%%%%%%%%%%%%%%%

%\documentclass[12pt,titlepage,onepage, letterpaper]{article}

\documentclass[12pt,titlepage,onepage, letterpaper]{report}


%
%-- specify related packages
%
\usepackage[T1]{fontenc}
\usepackage[USenglish]{babel}
\usepackage[nodayofweek,level]{datetime}
\usepackage[utf8]{inputenc}
\usepackage{subfig}
\usepackage{caption}
\usepackage{multirow}
\usepackage{rotating}
\usepackage{pdfpages}
\usepackage{textcomp}
\usepackage{mathtools}
\usepackage{eso-pic}
\usepackage{subfig}
\usepackage{float}
\floatstyle{plaintop}
\restylefloat{table}

\providecommand{\pgfsyspdfmark}[3]{}

\usepackage{array}
\newcolumntype{L}[1]{>{\raggedright\let\newline\\\arraybackslash\hspace{0pt}}m{#1}}
\newcolumntype{C}[1]{>{\centering\let\newline\\\arraybackslash\hspace{0pt}}m{#1}}
\newcolumntype{R}[1]{>{\raggedleft\let\newline\\\arraybackslash\hspace{0pt}}m{#1}}

%
% \usepackage[utf8x]{inputenc}
%

\usepackage{apacite}           %-- APA style citation 
                               %-- refer to http://www.ctan.org/tex-archive/biblio/bibtex/contrib/apacite/

%
%  \usepackage{ucs}
%


\usepackage{amsmath}           %-- American Math Society packages
\usepackage{amsfonts}
\usepackage{amssymb}


\usepackage{graphicx}          %-- graphicx package needed for including figures in JPG or PNG format
 
%
%\usepackage{graphics}          %-- graphics related package (this was commented out) use when image is in EPS format
%

\usepackage{verbatim}          %-- this package allows you to have multiple lines of comments by
                               %-- example:
                               %   \begin{comment}
                               %        ...your text here...
                               %   \end{comment}  

\usepackage{color}             %-- allows use of color with text
                               %-- example:  \textcolor{red}{This is the colored text in red.}

\usepackage{url}  %-- allows use of URLs example: \url{https:\ccs1.dlsu.edu.ph}

\usepackage{filecontents,lipsum}
\begin{filecontents*}{mystyle.sty}
\NeedsTeXFormat{LaTeX2e}
\ProvidesPackage{mystyle}
\newenvironment{acknowledgements}
  {\renewcommand{\abstractname}{Acknowledgements}% Abstract > Acknowledgements
   \begin{abstract}}
  {\end{abstract}
   \clearpage}
\endinput
\end{filecontents*}
\usepackage{mystyle}


%
%-- set margins,  you may need to edit this for your own printer
%
\topmargin 0.65in
\oddsidemargin 0.5in
\evensidemargin 0.5in

\voffset 0.0in
\hoffset 0.0in

\textwidth 5.75in
\textheight 7.0in


\parskip 1em
\parindent 0.25in

\bibliographystyle{apacite}            %-- use APA citation scheme

\hyphenation{ana-lysis know-ledge}     %-- LaTeX may not hyphenate correctly some words you use in your document
                                       %-- use \hyphenation to instruct LaTeX how to do it correctly, example above

\newcommand{\degree}{^{\circ}}         %-- use \newcommand to create your own "commands"
                                       %-- \newcommand works like the #define you learned in your COMPRO1 class

\newcommand{\etal}{et al.}


%\newcommand{\sinag}{\emph{Sinag}}
%\newcommand{\sinagtwo}{\emph{Sinag2}}

\newcommand{\figref}[1]{Figure \ref{#1}}
\newcommand{\appref}[1]{Appendix \ref{#1}}

%-- \newcommand{\Section}[1]{\section{#1}\setcounter{figure}{0}\setcounter{table}{0}}

%\newcommand{\shade}{\multicolumn{1}{|>{\columncolor[gray]{0.25}}c|}{}}
%\newcommand{\tableheader}[1]{\rowcolor{black}\color{white}{#1}}
%\newcommand{\cell}[2]{\multicolumn{1}{#1}{#2}}
%\newcommand{\definition}[2]{\textbf{\textit{#1}} --- #2}
%\newcommand{\itembit}[1]{\item \textbf{\textit{#1}}}
%\newcommand{\sgdef}[2]{\parbox[t][][t]{1.75in}{\textbf{#1}} \> \parbox[t][][t]{4.0in}{#2}\\\\}

%\newenvironment{sinagglossary}{\begin{flushleft}
%\begin{tabbing}
%\hspace{1.75in}\=\\}{\end{tabbing}\end{flushleft}}

\newcommand{\thestitle}[1]{{\Large \textsc{#1}}}


%---
%  \renewcommand{\thefigure}{\thesection.\arabic{figure}}
%  \renewcommand{\thetable}{\thesection.\arabic{table}}
%  \renewcommand{\contentsname}{Table of Contents}

\newcommand{\ts}{\textsuperscript}

\DeclarePairedDelimiter\abs{\lvert}{\rvert}


% \backgroundsetup{scale = 1, angle = 0, opacity = 0.2,
%    contents = {\includegraphics[width = \paperwidth,
%    height = \paperheight, keepaspectratio]
%    {dlsu-thesis-paper-letterhead.pdf}}}

\AddToShipoutPictureBG{\includegraphics[width=\paperwidth,height=\paperheight]{dlsu-thesis-paper-letterhead.pdf}}